\section{Discussion}

\textbf{Where to spend the parameters.} In this MoE model, the most effective place for
trainable parameters is the experts that already process most of the tokens, updated at full
rank. A small number of parameters works well in the attention, where every token passes,
but the same number spread thinly over all experts does not. Both observations point the
same way: what limits an adapter in a sparse model is how much of the computation it reaches.
Usage-based selection exploits this directly, which is why the original ESFT score works and
our specificity score does not. The result that the most used experts are not
domain-specific agrees with the high router entropy we measure: in a fine-grained model with
top-8 routing, domain knowledge is spread over many shared experts rather than isolated in a
few.

\textbf{Negative result for ESFT-Specific.} Ranking experts by $\log(p_\text{fin}/p_\text{gen})$
finds experts that are indeed more typical of finance, but they are rarely chosen, so their
updates affect few tokens. A score that combines both factors, e.g.\ usage weighted by
specificity, or the same selection with a larger budget, may work better; we did not test it.
Specificity might still be useful when catastrophic forgetting of the general domain matters,
because the selected experts carry only 5\% of the general-domain routing (vs.\ 21\% for ESFT),
but we did not measure forgetting.

\textbf{MoE vs.\ dense.} At the scale we can afford, the MoE model is not better than a dense
model with a similar number of active parameters: ESFT matches the dense model in mean
macro-F1 and accuracy but needs 60 times more trainable parameters, twice the memory and more than
twice the latency. The dense model is clearly better on Headlines, the MoE model on \fiqa{}. The
efficiency gap is partly due to our expert implementation (a loop instead of a fused grouped
GEMM) and would be smaller with optimised kernels.

\textbf{Limitations.}
\begin{itemize}
  \item One MoE model (1.3B) and one dense model; results for large MoE models with fewer,
  bigger experts may differ.
  \item No hyper-parameter search; learning rates follow common practice per method. The
  router may need a much lower learning rate or a regulariser to keep the routing close to
  the pre-trained one, and its failure here is a result for this setting only.
  \item The ESFT variants have about 47 times more trainable parameters than the LoRA methods,
  so the comparison between parameter \emph{locations} is budget-matched only for attention
  vs.\ expert LoRA.
  \item Three seeds; differences below about one standard deviation (e.g.\ ESFT vs.\ the dense
  model) are not significant. Headlines is the most seed-sensitive task, and the mean F1 of
  several methods is dominated by it.
  \item Short single-epoch training on 8.9K examples; we do not evaluate forgetting on general
  data or other financial tasks such as named-entity recognition or question answering.
\end{itemize}

\section{Conclusion}
We compared where to place trainable parameters when adapting a sparse MoE language model to
financial text classification. Updating the experts that the model already uses most (ESFT)
gives the best and the most stable result, comparable to LoRA on a dense model of similar
active size; attention LoRA is a strong cheap alternative, expert LoRA at the same budget is
weaker, and tuning only the router is harmful. Selecting domain-specific instead of frequently
used experts performs worse than random selection, and a routing analysis shows the reason:
the benefit of a tuned expert is proportional to the share of tokens routed to it. For
practitioners, this suggests choosing adapted experts by the amount of computation they
handle on the target data, not by how specialised they look.
